\begin{question}
\noindent $A$, $B$, $C$ and $D$ are points on the circumference of a circle, centre $O$.
$PA$ and $PD$ are tangents to the circle at $A$ and $D$ respectively.
Angle $ABC = 118^\circ$ and angle $AOD = 74^\circ$ (where $AOD$ is the reflex angle on the side containing $C$).

\begin{center}
\begin{tikzpicture}[scale=2.4]
    \coordinate (O) at (0,0);

    % Points on circle at chosen angles
    \coordinate (A) at (150:1);
    \coordinate (D) at (76:1);
    \coordinate (B) at (230:1);
    \coordinate (C) at (300:1);

    \draw[fill=lightgray!20] (O) circle (1cm);
    \fill[black] (O) circle (0.6pt);

    % Cyclic quadrilateral ABCD
    \draw (A) -- (B) -- (C) -- (D) -- cycle;

    % Radii to A and D
    \draw (O) -- (A);
    \draw (O) -- (D);

    % Tangent lines at A and D meeting at P
    \coordinate (TA) at ($(A)!3.2!-90:(O)$);
    \coordinate (TD) at ($(D)!3.2!90:(O)$);
    \path[name path=tanA] (A) -- (TA);
    \path[name path=tanD] (D) -- (TD);
    \path[name intersections={of=tanA and tanD, by=P}];
    \draw (A) -- (P) -- (D);

    % Right angle markers at tangent points
    \tkzMarkRightAngle[size=.15](O,A,P);
    \tkzMarkRightAngle[size=.15](P,D,O);

    % Labels
    \node[above left] at (A) {$A$};
    \node[above right] at (D) {$D$};
    \node[below left] at (B) {$B$};
    \node[below right] at (C) {$C$};
    \node[right] at (O) {$O$};
    \node[above] at (P) {$P$};

    \node at (0.15,-0.55) {$74^\circ$};
    \node at (-0.55,-0.15) {$118^\circ$};

\end{tikzpicture}
\end{center}

\noindent Show that angle $APD = 106^\circ$.

\noindent You must give a reason for each stage of your working, naming the circle theorem used.
\vspace{5cm}

\begin{flushright}
\answermarks{4}
\end{flushright}
\end{question}
